\documentclass[12pt]{article}
%%%%%%%%%%%%%%%%%%%%%%%%%%%%%%%%%%%%%%%%%%%%%%%%%%%%%%%%%%%%%%%%%%%%%%%%%%%%%%%%%%%%%%%%%%%%%%%%%%%%%%%%%%%%%%%%%%%%%%%%%%%%%%%%%%%%%%%%%%%%%%%%%%%%%%%%%%%%%%%%%%%%%%%%%%%%%%%%%%%%%%%%%%%%%%%%%%%%%%%%%%%%%%%%%%%%%%%%%%%%%%%%%%%%%%%%%%%%%%%%%%%%%%%%%%%%
\usepackage{amsfonts}
\usepackage{eurosym}
\usepackage{geometry}
\usepackage{amsmath,amsthm,amssymb}
\usepackage{ulem} 
\usepackage{graphicx}
\usepackage{comment}
%\usepackage[sort,comma]{natbib}
\usepackage[utf8]{inputenc}
\usepackage{setspace}
\usepackage[backend=biber, style = apa]{biblatex}
\usepackage{placeins} % to separate sections

\usepackage{adjustbox}
\usepackage{array}
\usepackage{multirow}
\usepackage{graphicx}
\usepackage{subcaption}
\usepackage{pifont}
\usepackage{amssymb}
\usepackage{comment}
\usepackage[hang, flushmargin, bottom]{footmisc}
\usepackage{footnotebackref}
\usepackage{xcolor}
\usepackage{hyperref}
\usepackage{booktabs}
\usepackage{pifont}
\usepackage{caption}
\usepackage{float}
\usepackage{todonotes}
\setcounter{MaxMatrixCols}{10}


%\setlength{\bibsep}{0.3pt}
\setlength{\textfloatsep}{5pt}
\hypersetup{breaklinks=true,hypertexnames=false,colorlinks=true,citecolor = teal}
\captionsetup{font=normalsize}
\newcommand{\cmark}{\ding{51}}
\def\sym#1{\ifmmode^{#1}\else\(^{#1}\)\fi}
\renewcommand{\thetable}{\Roman{table}}
\geometry{verbose,tmargin=.9in,bmargin=1in,lmargin=.8in,rmargin=.8in,nomarginpar}
\makeatletter
\DeclareTextSymbolDefault{\textquotedbl}{T1}
\theoremstyle{plain}
\newtheorem{thm}{\protect\theoremname}
\theoremstyle{plain}
\newtheorem{prop}[thm]{\protect\propositionname}
\theoremstyle{definition}  % Add this line
\newtheorem{definition}[thm]{Definition}  % Add this line
\theoremstyle{remark}  % Add this line
\newtheorem{remark}[thm]{Remark}  % Add this line
\providecommand{\propositionname}{Proposition}
\newtheorem{proposition}{Proposition}

\providecommand{\theoremname}{Theorem}
\makeatother
\newtheorem{ass}[thm]{Assumption}
\newtheorem{lemma}[thm]{Lemma}
\newtheorem{corollary}[thm]{Corollary}
% \input{tcilatex}
\usepackage{tikz}
\usetikzlibrary{shapes.geometric, arrows, positioning}


\addbibresource{../references.bib}
\begin{document}

\section*{Model 6 (v3): Identification and Estimation with Heterogeneous Covariates}

This document develops the identification theory for the structural model with observable heterogeneity (Model~6, \texttt{M6\_heterogeneity\_information.tex}) when the econometrician does not observe---or only partially observes---the public covariates~$\mathbf{x}$ that both banks use.

Model~5 (\texttt{M5\_convert4intostructural.tex}) establishes nonparametric identification when there is no observable heterogeneity: bank~2 plays a single rate, and three data objects $(r_2^*, \sigma^*(\cdot), D_i)$ suffice to recover $(\alpha, \lambda, F)$.  Model~6 introduces public covariates~$\mathbf{x}$ observed by both banks, which generates heterogeneity in bank~2's rate across borrowers.  The central question is: \emph{what can the econometrician recover when~$\mathbf{x}$ is partially or fully unobserved?}

Section~\ref{sec:recap} recaps the model. Section~\ref{sec:suffstat} establishes the sufficient statistic result.  Section~\ref{sec:nonpar_id} characterizes precisely what is and is not identified nonparametrically.  Section~\ref{sec:parametric_id} provides formal identification results under parametric assumptions, stating exact conditions.  Section~\ref{sec:likelihood} derives the likelihood.  Section~\ref{sec:simulation} describes the Monte Carlo validation.


%===========================================================================
%  SECTION 1: MODEL RECAP
%===========================================================================

\section{Model Recap}\label{sec:recap}

Two banks compete for a borrower. Bank~1 (informed/incumbent) observes both public covariates $\mathbf{x}$ and private covariates $\mathbf{x}'$, and therefore knows the borrower's default probability $h = h(\mathbf{x}, \mathbf{x}')$. Bank~2 (uninformed/entrant) observes only $\mathbf{x}$ and faces the conditional distribution $F(h \mid \mathbf{x})$. Demand is logit:
\begin{align}\label{eq:share}
    s(r_1, r_2) = \frac{1}{1 + \exp\!\big(\alpha(r_1 - r_2) - \lambda\big)}
\end{align}
In equilibrium:
\begin{align}
    \sigma^*(h;\mathbf{x}) &= \frac{1}{1-h} + \frac{1}{\alpha\big(1-s(\sigma^*(h;\mathbf{x}), r_2^*(\mathbf{x}))\big)} \label{eq:sigma}\\[4pt]
    r_2^*(\mathbf{x}) &= \frac{1}{1 - \tilde{h}(\mathbf{x})} + \frac{1}{\alpha}\cdot\frac{\int (1-s)(1-h)\, dF(\cdot|\mathbf{x})}{\int s\,(1-s)(1-h)\, dF(\cdot|\mathbf{x})} \label{eq:r2}
\end{align}
The structural parameters are $\theta = (\alpha, \lambda, F(\cdot|\mathbf{x}))$.

\medskip\noindent
\textbf{Data.} For each borrower $i = 1, \ldots, N$, the econometrician observes:
\begin{center}
\renewcommand{\arraystretch}{1.3}
\begin{tabular}{lll}
\toprule
\textbf{Variable} & \textbf{Description} & \textbf{Values} \\
\midrule
$r_i^{\text{win}}$ & Interest rate charged by the chosen bank & $\mathbb{R}_+$ \\
$W_i$ & Identity of the chosen bank & $\{1, 2\}$ \\
$D_i$ & Default indicator & $\{0, 1\}$ \\
$\mathbf{z}_i$ & Observed public covariates (possibly empty) & $\mathbb{R}^{d_z}$ \\
\bottomrule
\end{tabular}
\end{center}
We partition $\mathbf{x} = (\mathbf{z}, \mathbf{w})$, where $\mathbf{z}$ is observed by the econometrician and $\mathbf{w}$ is not. The analysis covers all cases from $\mathbf{z} = \mathbf{x}$ (full observability, nonparametric identification as in M5) to $\mathbf{z} = \emptyset$ (no covariates, most demanding).

\medskip\noindent
\textbf{Key data asymmetry:}
\begin{itemize}
    \item When $W_i = 2$: we observe $r_i^{\text{win}} = r_2^*(\mathbf{x}_i)$, directly revealing the entrant's rate for that borrower.
    \item When $W_i = 1$: we observe $r_i^{\text{win}} = \sigma^*(h_i; r_2^*(\mathbf{x}_i))$, which depends on two latent variables: the type $h_i$ and the entrant rate $r_2^*(\mathbf{x}_i)$.
\end{itemize}


%===========================================================================
%  SECTION 2: THE SUFFICIENT STATISTIC
%===========================================================================

\newpage
\section{The Sufficient Statistic Result}\label{sec:suffstat}

\begin{prop}[Sufficient statistic]\label{prop:suffstat}
In the stage-2 equilibrium, the public covariates $\mathbf{x}$ affect bank~1's pricing rule $\sigma^*$ and the market shares $s$ only through the entrant's rate $r_2^*(\mathbf{x})$. Formally, for any two covariate vectors $\mathbf{x}_a, \mathbf{x}_b \in \mathcal{X}$ with $r_2^*(\mathbf{x}_a) = r_2^*(\mathbf{x}_b) = r_2$:
\begin{align}\label{eq:suffstat}
    \sigma^*(h; \mathbf{x}_a) = \sigma^*(h; \mathbf{x}_b) = \sigma^*(h; r_2), \qquad s(\sigma^*(h; \mathbf{x}_a), r_2) = s(\sigma^*(h; \mathbf{x}_b), r_2)
\end{align}
for all $h$.
\end{prop}

\begin{proof}
Immediate from the bank~1 FOC \eqref{eq:sigma}: the right side depends on $\mathbf{x}$ only through $r_2^*(\mathbf{x})$. Since bank~1's best response is the unique solution to this fixed-point equation, and the equation is identical for all $\mathbf{x}$ with the same $r_2^*$, the solution $\sigma^*(h; \mathbf{x})$ is the same.
\end{proof}

\begin{remark}[What $r_2^*$ does not capture]\label{rem:suffstat_limit}
The sufficient statistic result applies to bank~1's strategy and market shares, but \emph{not} to the full equilibrium. Bank~2's FOC \eqref{eq:r2} depends on $F(h|\mathbf{x})$, which is not a function of $r_2^*$ alone---two different~$\mathbf{x}$ values can lead to different $F(h|\mathbf{x})$ but the same optimal $r_2^*$.
\end{remark}

\begin{definition}[Type-recovery function]\label{def:tau}
For fixed $r_2$, define $\tau(r_1; r_2)$ as the inverse of $\sigma^*(\cdot\,; r_2)$:
\begin{align}\label{eq:tau_def}
    \tau\!\big(\sigma^*(h; r_2);\, r_2\big) = h \qquad \text{for all } h
\end{align}
Since $\sigma^*(\cdot; r_2)$ is strictly increasing in $h$, $\tau(\cdot; r_2)$ is well-defined and strictly increasing in $r_1$.  The Jacobian is:
\begin{align}\label{eq:tau_prime}
    \tau'(r_1; r_2) = \frac{1}{\sigma^*_h\!\big(\tau(r_1;r_2);\, r_2\big)} > 0
\end{align}
\end{definition}


%===========================================================================
%  SECTION 3: NONPARAMETRIC IDENTIFICATION
%===========================================================================

\newpage
\section{What Is and Is Not Identified Nonparametrically}\label{sec:nonpar_id}

We now characterize precisely what can be recovered from the data without parametric assumptions, and where nonparametric identification breaks down. This makes explicit what the parametric assumptions in Section~\ref{sec:parametric_id} are doing.

\subsection{Benchmark: Full Observability ($\mathbf{z} = \mathbf{x}$)}\label{sec:benchmark}

When the econometrician observes all public covariates, conditioning on $\mathbf{x}$ reduces the model to M5 within each $\mathbf{x}$-stratum. The M5 identification result applies market-by-market:

\begin{enumerate}
    \item $r_2^*(\mathbf{x})$ is identified as the common winning rate among bank-2 winners with covariates $\mathbf{x}$.
    \item $\sigma^*(\cdot; \mathbf{x})$ is identified via type recovery: $E[D_i \mid r_i^{\text{win}} = r, W_i = 1, \mathbf{x}_i = \mathbf{x}] = \tau(r; r_2^*(\mathbf{x}))$.
    \item $(\alpha, \lambda)$ are identified from the linear restriction $\ln(\alpha\,m -1) = \lambda - \alpha\Delta$ within any single $\mathbf{x}$-stratum.  Multiple strata provide overidentification.
    \item $F(h|\mathbf{x})$ is identified by inverting the selection-weighted density of types among bank-1 winners.
\end{enumerate}

\noindent This is the nonparametric ideal. We now analyze what survives when $\mathbf{x}$ is partially or fully unobserved.


\subsection{What IS Identified Without $\mathbf{x}$}\label{sec:what_is}

Even without observing $\mathbf{x}$, several objects are nonparametrically identified from the data. Understanding these is important because they constrain any parametric model.

\begin{prop}[Identified objects without covariates]\label{prop:identified_objects}
From the joint distribution of $(r_i^{\text{win}}, W_i, D_i)$ alone, the following are nonparametrically identified:
\begin{enumerate}
    \item[(a)] The distribution of $r_2^*(\mathbf{x})$ among bank-2 winners (the selected distribution).
    \item[(b)] The conditional default rate $E[D_i \mid r_i^{\text{win}}, W_i = 2]$ as a function of the winning rate.
    \item[(c)] The distribution of $\sigma^*(h_i; r_2^*(\mathbf{x}_i))$ among bank-1 winners.
    \item[(d)] The conditional default rate $E[D_i \mid r_i^{\text{win}}, W_i = 1]$ as a function of the winning rate.  However, this does \emph{not} recover the borrower's type~$h$.  The reason is that bank~1's rate depends on both the borrower's type and the competitive environment: $r_1 = \sigma^*(h; r_2^*(\mathbf{x}))$.  Two borrowers with the \emph{same} type~$h$ but different covariates~$\mathbf{x}$ face different entrant rates $r_2^*(\mathbf{x})$, so bank~1 charges them different rates~$r_1$.  Conversely, the same~$r_1$ can arise from different $(h, \mathbf{x})$ combinations: a safe borrower whose covariates~$\mathbf{x}$ imply a competitive entrant (low~$r_2^*$), or a riskier borrower whose covariates imply a less competitive entrant (high~$r_2^*$).  Since the econometrician does not observe~$\mathbf{x}$, she cannot determine which competitive environment generated the observed~$r_1$, and therefore cannot back out~$h$.  The conditional default rate $E[D_i \mid r_1, W_i\!=\!1]$ is thus a weighted average over the different types that produce rate~$r_1$, not the type itself.
    \item[(e)] The aggregate switching probability $\Pr(W_i = 2)$.
\end{enumerate}
\end{prop}

\begin{proof}
These are all features of the observable joint distribution $(r_i^{\text{win}}, W_i, D_i)$ and are identified directly from the data without any model structure.
\end{proof}

Let us examine each and its informational content.

\medskip\noindent
\textbf{(a) Distribution of $r_2^*$ among switchers.}  Among bank-2 winners, $r_i^{\text{win}} = r_2^*(\mathbf{x}_i)$, so the empirical distribution of winning rates directly reveals the \emph{selected} distribution.  However, this is \emph{not} the population distribution of $r_2^*$: the density observed is
\begin{align}\label{eq:g2_selected}
    g_{r_2 | W=2}(r_2) = \frac{g(r_2) \cdot \Pr(W=2 \mid r_2)}{\Pr(W=2)}
\end{align}
where $g(r_2)$ is the population density and $\Pr(W=2 \mid r_2) = \int (1-s(\sigma^*(h; r_2), r_2))\, dF(h|r_2)$ is the aggregate switching probability conditional on $r_2$.  Lower $r_2$ (more competitive entrant) means more switching, so the observed distribution is shifted downward relative to the population.  Recovering $g(r_2)$ from $g_{r_2|W=2}(r_2)$ requires knowing $\Pr(W=2|r_2)$, which depends on $(\alpha, \lambda, F)$.

\medskip\noindent
\textbf{(b) Default rates among bank-2 winners.}  For each value of the winning rate $r_2$:
\begin{align}\label{eq:d2_mixture}
    E[D_i \mid r_i^{\text{win}} = r_2, W_i = 2] = \frac{\int h\,(1-s(\sigma^*(h;r_2),r_2))\, dF(h|r_2)}{\int (1-s(\sigma^*(h;r_2),r_2))\, dF(h|r_2)}
\end{align}
This is the \emph{adverse-selection-weighted} mean type: borrowers who switch to bank~2 are disproportionately those with high~$h$ (risky types), because they face a higher $\sigma^*(h)$ that pushes them away from bank~1.  Formally, $1-s(\sigma^*(h),r_2)$ is increasing in~$h$ (since $\sigma^*$ is increasing and $s_1 < 0$), so $E[D|r_2, W\!=\!2] > E[h|r_2]$.  This object is directly estimable by regressing~$D_i$ on~$r_i^{\text{win}}$ among bank-2 winners.

\medskip\noindent
\textbf{(d) The mixture problem for bank-1 winners.}  Among bank-1 winners, the winning rate $r_1 = \sigma^*(h_i; r_2^*(\mathbf{x}_i))$ confounds two latent variables: $h_i$ and $r_2^*(\mathbf{x}_i)$.  The same $r_1$ can arise from different $(h, r_2)$ combinations:
\begin{itemize}
    \item Low $h$ (good type) + low $r_2$ (competitive entrant): bank~1 offers a low rate against tough competition.
    \item Higher $h$ (riskier) + higher $r_2$ (less competitive entrant): bank~1 charges more due to both higher cost and more market power.
\end{itemize}
The conditional default rate is a mixture:
\begin{align}\label{eq:mixture}
    E[D_i \mid r_i^{\text{win}} = r_1, W_i = 1] = \frac{\int_{\mathcal{R}_2(r_1)} \tau(r_1; r_2) \cdot s(r_1, r_2) \cdot f(\tau(r_1;r_2)|r_2) \cdot \tau'(r_1;r_2) \cdot g(r_2)\, dr_2}{\int_{\mathcal{R}_2(r_1)} s(r_1, r_2) \cdot f(\tau(r_1;r_2)|r_2) \cdot \tau'(r_1;r_2) \cdot g(r_2)\, dr_2}
\end{align}
This \emph{cannot} be inverted to recover $h_i$ for any individual borrower without additional structure.


\subsection{What Fails Without $\mathbf{x}$: The Benchmark Argument}\label{sec:what_fails}

Without observing~$\mathbf{x}$, the four-step nonparametric identification argument from the benchmark case (Section~\ref{sec:benchmark}) breaks down:
\begin{enumerate}
    \item[(i)] \textbf{Step 1 (Recovery of $r_2^*$):} Bank~2's rate is no longer a single constant. It varies across borrowers through $\mathbf{x}$, generating a distribution rather than a point mass among bank-2 winners.
    
    \item[(ii)] \textbf{Step 2 (Type recovery):} The key identity $E[D_i \mid r_1, W_i=1] = h_i$ fails.  The source of the problem is heterogeneity in the competitive environment through~$\mathbf{x}$.  Borrowers with covariates~$\mathbf{x}$ that make the entrant very competitive (low~$r_2^*(\mathbf{x})$) face a very different pricing environment from borrowers whose covariates make the entrant less competitive (high~$r_2^*(\mathbf{x})$).  Bank~1's rate $r_1 = \sigma^*(h; r_2^*(\mathbf{x}))$ depends on both the type and the competitive environment, and $\sigma^*$ is increasing in both arguments.  So the same~$r_1$ can be generated by a low-risk borrower in a competitive market (low~$h$, low~$r_2^*$) or a higher-risk borrower in a less competitive market (high~$h$, high~$r_2^*$).  Without observing~$\mathbf{x}$, the econometrician cannot determine which competitive environment generated the observed rate, so the conditional expectation $E[D_i \mid r_1, W_i\!=\!1]$ averages over these different types rather than recovering a single~$h$.
    
    \item[(iii)] \textbf{Step 3 (Identifying restriction):} The linear restriction $\ln(\alpha m - 1) = \lambda - \alpha\Delta$ requires knowing both $m(h) = \sigma^*(h) - 1/(1-h)$ and $\Delta(h) = \sigma^*(h) - r_2^*$ for individual borrowers. Without type recovery, neither $m$ nor $\Delta$ can be computed.
    
    \item[(iv)] \textbf{Step 4 (Distribution recovery):} Recovering $F(h)$ from the selection-weighted density $f_{h|W=1}(h)$ requires knowing $s(h)$, which requires $(\alpha, \lambda)$ from Step~3.
\end{enumerate}

However, showing that a particular identification strategy fails does not establish that the model is unidentified---perhaps a different argument could work.  We now prove that no such argument exists.


\subsection{Nonparametric Non-Identification}\label{sec:nonid_formal}

\begin{prop}[Nonparametric non-identification]\label{prop:nonid}
Without parametric restrictions on $F(h|\mathbf{x})$ and the distribution of~$\mathbf{x}$, the structural parameters $(\alpha, \lambda)$ are not identified.  Specifically: for any parameter vector $(\alpha_0, \lambda_0, F_0, G_0)$ generating observable distribution~$\mathcal{D}$, and for any $(\alpha_1, \lambda_1)$ in a neighborhood of~$(\alpha_0, \lambda_0)$, there exist nonparametric distributions $F_1$ and~$G_1$ such that $(\alpha_1, \lambda_1, F_1, G_1)$ also generates~$\mathcal{D}$.
\end{prop}

\begin{proof}
We construct an explicit observationally equivalent model.  Fix $(\alpha_1, \lambda_1) \neq (\alpha_0, \lambda_0)$ with $\alpha_1 > 0$.

\medskip\noindent
\textbf{Step 1: New equilibrium.}  Under $(\alpha_1, \lambda_1)$, bank~1's FOC defines a new equilibrium strategy $\sigma_1^*(h; r_2)$ and market shares $s_1(r_1, r_2)$ for each~$r_2$.  These are uniquely determined (for given $F$ and $r_2$, the bank-1 FOC has a unique solution under standard regularity).

\medskip\noindent
\textbf{Step 2: Matching bank-2 winner defaults.}  The data identify $\delta(r_2) \equiv E[D_i \mid r_2, W_i = 2]$ at each~$r_2$.  Under the alternative model:
\[
    \delta(r_2) = \frac{\int h\,(1 - s_1(\sigma_1^*(h; r_2), r_2))\, dF_1(h|r_2)}{\int (1 - s_1(\sigma_1^*(h; r_2), r_2))\, dF_1(h|r_2)}
\]
This is a single moment restriction on the nonparametric distribution~$F_1(\cdot|r_2)$ at each~$r_2$.  Since $F_1(\cdot|r_2)$ is a probability distribution on $[\underline{h}, \bar{h}]$---an infinite-dimensional object---a single moment does not pin it down.  We can choose $F_1(\cdot|r_2)$ to satisfy $\delta(r_2)$ at each~$r_2$, with many degrees of freedom remaining.

\medskip\noindent
\textbf{Step 3: Matching the $r_2$ distribution.}  Given $(\alpha_1, \lambda_1, F_1)$, the selection probability $\Pr_1(W\!=\!2|r_2)$ is determined.  Set:
\[
    g_1(r_2) \propto \frac{g_{r_2|W=2}(r_2)}{\Pr_1(W\!=\!2|r_2)}
\]
where $g_{r_2|W=2}$ is the observed density of $r_2$ among bank-2 winners.  This ensures the alternative model matches the observed distribution of entrant rates among switchers.

\medskip\noindent
\textbf{Step 4: Matching bank-1 winner density.}  The density of $(r_1, D)$ among bank-1 winners is a mixture:
\[
    p(r_1, d) = \int f_1(\tau_1(r_1; r_2) \mid r_2)\, \tau_1'(r_1; r_2)\, s_1(r_1, r_2)\, h^d(1\!-\!h)^{1-d}\, g_1(r_2)\, dr_2
\]
Here $\tau_1(\cdot; r_2)$ is the type-recovery function under the alternative equilibrium.  At each~$r_2$, the integrand involves $f_1(\cdot|r_2)$, which has the remaining degrees of freedom from Step~2.  The system of equations (one per $r_1$ value) constrains the marginal contributions at each~$r_2$ through a Fredholm integral equation of the first kind.  Since $f_1(\cdot|r_2)$ is unrestricted (beyond the single moment matched in Step~2), this system generically admits solutions: the operator maps an infinite-dimensional input to a one-dimensional output at each~$r_1$.

\medskip\noindent
\textbf{Conclusion.}  For any $(\alpha_1, \lambda_1)$ near the truth, we can find nonparametric $(F_1, G_1)$ matching all observable moments:
\begin{itemize}
    \item The distribution of $r_2$ among bank-2 winners (Step~3),
    \item The conditional default rates among bank-2 winners (Step~2),
    \item The joint density of $(r_1, D)$ among bank-1 winners (Step~4),
    \item The aggregate switching rate (absorbed by the normalization of $G_1$ in Step~3).
\end{itemize}
The fundamental reason is a \emph{dimension mismatch}: the data provide functional restrictions (a continuum indexed by $r_1$ or $r_2$), but the nonparametric unknowns $F_1(h|r_2)$ constitute a family of distributions---one per~$r_2$---which is a higher-dimensional object.  Two scalar parameters $(\alpha, \lambda)$ cannot be pinned down when the nuisance functions absorb all identifying variation.
\end{proof}

\begin{remark}[Intuition]
The non-identification has a simple economic interpretation.  Suppose the true $\alpha$ (price sensitivity) is moderate.  An alternative model with higher~$\alpha_1$ would predict stronger adverse selection (riskier borrowers switch more), changing the predicted default rates.  But this can be offset by choosing $F_1(h|r_2)$ with a different shape---for instance, a type distribution with less dispersion, so that even with stronger selection, the default rates among switchers match the data.  With nonparametric $F$, there is always enough flexibility to absorb any change in $(\alpha, \lambda)$.
\end{remark}


\subsection{The Role of Default Data in Narrowing Identification}\label{sec:defaults_nonpar}

Default outcomes $D_i$ provide additional identifying power beyond rates and winner identity. Two mechanisms:

\begin{enumerate}
    \item \textbf{Among bank-2 winners:} The function $E[D_i \mid r_i^{\text{win}} = r_2, W_i = 2]$ is identified. Since this is a weighted average of types who switch at entrant rate $r_2$, it traces out how the adverse-selection-weighted mean type varies with $r_2$.  A structural model must match this entire curve, not just its level---this is a functional restriction, not a scalar one.
    
    \item \textbf{Among bank-1 winners:} Although $E[D_i \mid r_1, W_i = 1]$ is a mixture (not a clean type recovery), its \emph{shape} is informative.  For instance:
    \begin{itemize}
      \item If the default rate rises steeply with $r_1$, this is consistent with a narrow range of $r_2^*$ (the mixture is concentrated, close to a clean type recovery).
      \item If the default rate is flat or non-monotone, this indicates substantial variation in $r_2^*$ (the mixture blurs the $h$--$r_1$ mapping).
    \end{itemize}
    The \emph{variance} $\text{Var}(D_i \mid r_1, W_i = 1)$ is also informative: if limited variation in $r_2$ only one type produces each $r_1$, then $\text{Var}(D_i \mid r_1, W_i=1) \approx h(1-h)$ (Bernoulli variance). Excess variance beyond $h(1-h)$ indicates mixture heterogeneity in $h$ at that rate.
\end{enumerate}

\subsection{Summary: Identification Landscape}\label{sec:landscape}

\begin{center}
\renewcommand{\arraystretch}{1.5}
\begin{tabular}{p{4cm}p{3.5cm}p{5.5cm}}
\toprule
\textbf{Object} & \textbf{Status} & \textbf{What's missing} \\
\midrule
$\alpha, \lambda$ & Not point-identified & Require clean type recovery (Step 2), which fails without $r_2$ \\[2pt]
$g_{r_2|W=2}(\cdot)$ & Point-identified & Directly from bank-2 winning rates \\[2pt]
$g(r_2)$ (population) & Not point-identified & Need $\Pr(W\!=\!2|r_2)$, which requires $(\alpha,\lambda,F)$ \\[2pt]
$E[D|r_2, W\!=\!2]$ & Point-identified & Identified as a function of $r_2$ \\[2pt]
$F(h|\mathbf{x})$ & Not point-identified & Requires $(\alpha,\lambda)$ and clean type recovery \\[2pt]
$\Pr(W\!=\!2)$ & Point-identified & Sample proportion \\
\bottomrule
\end{tabular}
\end{center}

\noindent
The upshot: without parametric restrictions, we identify the \emph{marginals} and \emph{conditional moments} of the observable distribution, but cannot disentangle the latent variables $(h_i, r_2^*(\mathbf{x}_i))$ for bank-1 winners.  This motivates the parametric approach in the next section.


%===========================================================================
%  SECTION 4: PARAMETRIC IDENTIFICATION
%===========================================================================

\newpage
\section{Parametric Identification}\label{sec:parametric_id}

We now show that imposing parametric structure on two specific objects---the type distribution and the mixing distribution---restores identification, and we characterize the conditions under which this holds.

\subsection{Parametric Specification}\label{sec:spec}

\begin{ass}[Parametric specification]\label{ass:parametric}
\begin{enumerate}
    \item[(i)] \textbf{Type distribution.} $h \mid \mathbf{x} \sim F(h;\, \theta_F(\mathbf{x}))$, where $F$ belongs to a known parametric family and $\theta_F(\mathbf{x})$ varies with~$\mathbf{x}$ through a finite-dimensional parameterization.
    
    \textbf{Example:} $h \mid x \sim \text{Beta}(a,\, b(x))$ on $[h_{\min}, h_{\max}]$, with $b(x) = b_0 + b_1 x$, giving $\theta_F = (a, b_0, b_1, h_{\min}, h_{\max})$.
    
    \item[(ii)] \textbf{Covariate distribution.} $\mathbf{x} \sim P(\mathbf{x};\, \psi)$ belongs to a known parametric family.
    
    \textbf{Example:} $x \sim \text{LogNormal}(\mu_x, \sigma_x^2)$, giving $\psi = (\mu_x, \sigma_x)$.
\end{enumerate}
When the econometrician observes some covariates $\mathbf{z} \subseteq \mathbf{x}$, the conditioning extends: $\theta_F(\mathbf{x}, \mathbf{z})$ and $P(\mathbf{x} | \mathbf{z}; \psi)$.
\end{ass}

The full parameter vector is $\theta = (\alpha, \lambda, \theta_F, \psi)$.

\begin{remark}[Reduced-form objects induced by the primitives]\label{rem:induced}
The primitives $(F(h|\mathbf{x}; \theta_F),\; P(\mathbf{x}; \psi))$ determine the equilibrium objects that enter the likelihood:
\begin{itemize}
    \item \textbf{Entrant rate.}  For each $\mathbf{x}$, bank~2's equilibrium rate $r_2^*(\mathbf{x})$ is determined by the FOC~\eqref{eq:r2}, which depends on $(\alpha, \lambda, F(\cdot|\mathbf{x}))$.  This defines a mapping $\mathbf{x} \mapsto r_2^*(\mathbf{x}; \theta)$.
    \item \textbf{Mixing distribution.} The population distribution of entrant rates $G(r_2; \theta)$ is the pushforward of $P(\mathbf{x}; \psi)$ through $r_2^*(\cdot; \theta)$.  It is \emph{not} a free object but is determined by the primitives through the equilibrium.
    \item \textbf{Conditional type distribution.}  By the sufficient statistic result (Proposition~\ref{prop:suffstat}), conditionally on $r_2$, the type distribution $F(h|r_2; \theta)$ is determined by averaging $F(h|\mathbf{x}; \theta_F)$ over all~$\mathbf{x}$ that map to the same~$r_2$.
\end{itemize}
The likelihood (Section~\ref{sec:likelihood}) is written in terms of these derived objects $G(r_2; \theta)$ and $F(h|r_2; \theta)$, because the sufficient statistic result ensures that~$\mathbf{x}$ enters the likelihood only through~$r_2^*(\mathbf{x})$.
\end{remark}

\begin{remark}[Why two parametric choices are needed]\label{rem:why_two}
\begin{itemize}
    \item \textbf{Assumption (i)} constrains the type distribution.  Proposition~\ref{prop:nonid} shows that without it, for any candidate $(\alpha, \lambda)$, there exist type distributions that rationalize the observed data.  The parametric family restricts the degree of freedom in $F$, enabling the data to discipline $(\alpha, \lambda)$.
    \item \textbf{Assumption (ii)} characterizes the population distribution of~$\mathbf{x}$, which determines the equilibrium distribution of competitive conditions.  While $r_2$ is observed among bank-2 winners, those observations are \emph{selected} (borrowers who switched), and recovering the population distribution requires knowing $\Pr(W\!=\!2|r_2)$---which depends on all structural parameters.  The parametric family anchors the covariate distribution.
\end{itemize}
Importantly, neither assumption alone suffices.  The identification works through the \emph{interaction}: the parametric $F(h|\mathbf{x})$ restricts how defaults relate to rates for each competitive environment, and the parametric $P(\mathbf{x})$ restricts which competitive environments arise and how frequently.
\end{remark}

\subsection{Restriction Counting: Why Parametric Models Are Generically Identified}\label{sec:formal_id}

Before analyzing the specific Beta specification, we establish a general result: any well-specified finite-dimensional parametric model satisfying Assumption~\ref{ass:parametric} is \emph{generically} identified---meaning that identification fails only on a measure-zero set in the parameter space.  This follows from a restriction-counting argument combined with the structure of the data.

\medskip\noindent
\textbf{Available restrictions.}  The data provide two families of functional restrictions:

\begin{enumerate}
    \item \textbf{From bank-2 winners:} For each $r_2$ in the support of bank-2 winners' rates, the model predicts the conditional default rate:
    \begin{align}\label{eq:d2_restriction}
        E[D|r_2, W\!=\!2] = \frac{\int h\,(1 - s(\sigma^*(h; r_2), r_2))\, dF(h;\theta_F(r_2))}{\int (1 - s(\sigma^*(h; r_2), r_2))\, dF(h;\theta_F(r_2))}
    \end{align}
    The left side is a known function of $r_2$ (estimable from data).  The right side depends on the structural parameters $(\alpha, \lambda, \theta_F)$.  If the support of $r_2$ has non-empty interior, this is a \emph{continuum} of restrictions---infinitely many equations constraining finitely many parameters.
    
    \item \textbf{From bank-1 winners:} The density of $(r_1, D)$ among bank-1 winners is:
    \begin{align}\label{eq:b1_density}
        p(r_1, d, W\!=\!1 \mid \theta) = \int f(\tau(r_1;r_2) \mid r_2; \theta) \cdot \tau'(r_1;r_2) \cdot s(r_1, r_2) \cdot h^d(1\!-\!h)^{1-d} \cdot g(r_2;\theta)\, dr_2
    \end{align}
    where $f(\cdot|r_2;\theta)$ and $g(r_2;\theta)$ are the induced conditional type distribution and mixing density (Remark~\ref{rem:induced}).  This provides additional functional restrictions involving all parameters $(\alpha, \lambda, \theta_F, \psi)$.
    
    \item \textbf{Selection-distorted mixing distribution:} The density of $r_2$ among bank-2 winners is:
    \begin{align}\label{eq:g2_selection}
        g_{r_2|W=2}(r_2) = \frac{g(r_2;\theta) \cdot \Pr(W\!=\!2|r_2;\, \alpha,\lambda,\theta_F)}{\Pr(W\!=\!2)}
    \end{align}
    which constrains $\psi$ jointly with all other parameters.
    
    \item \textbf{Aggregate switching rate:} $\Pr(W\!=\!2) = \int\!\int (1-s(\sigma^*(h;r_2),r_2))\, dF(h|r_2;\theta)\, dG(r_2;\theta)$ provides a scalar restriction on $(\alpha,\lambda,\theta_F,\psi)$.
\end{enumerate}

\medskip\noindent
\textbf{Dimension count.}  The parameter vector $\theta = (\alpha, \lambda, \theta_F, \psi)$ is finite-dimensional: $\dim(\theta) = 2 + \dim(\theta_F) + \dim(\psi)$.  For the Beta specification in the Monte Carlo ($\theta_F = (a, b_0, b_1)$, $\psi = (\mu_x, \sigma_x)$), $\dim(\theta) = 7$.  The data provide:
\begin{itemize}
    \item A continuum of restrictions from $E[D|r_2, W\!=\!2]$ (one per $r_2$ value in the support),
    \item A continuum from the bank-1 density $p(r_1, d, W\!=\!1)$,
    \item A functional restriction from $g_{r_2|W=2}(\cdot)$,
    \item A scalar from $\Pr(W\!=\!2)$.
\end{itemize}
The system is massively overdetermined: infinitely many independent restrictions for finitely many parameters.

\begin{thm}[Generic identification]\label{thm:generic_id}
Suppose the parametric families $F(h|\mathbf{x};\theta_F)$ and $P(\mathbf{x};\psi)$ are smooth in their parameters, and the support of $r_2^*$ among bank-2 winners has non-empty interior.  Then the set of parameter values $\theta = (\alpha, \lambda, \theta_F, \psi)$ at which the model is \emph{not} identified (i.e., there exists $\theta' \neq \theta$ generating the same distribution of observables) has Lebesgue measure zero in the parameter space.
\end{thm}

\begin{proof}[Proof sketch]
The functional restrictions above define a system of equations $\mathcal{M}(\theta) = \mathcal{D}$, where $\mathcal{M}(\theta)$ is the model-predicted functional (mapping parameters to the joint distribution of observables) and $\mathcal{D}$ is the data distribution.  This is a map from $\mathbb{R}^{d_\theta}$ to an infinite-dimensional function space.  By a standard transversality argument (cf.\ Rothenberg, 1971, \emph{Econometrica}, Theorem~1), if the Jacobian $\partial \mathcal{M}/\partial \theta$ has full column rank at $\theta_0$ (the ``local identification condition''), then $\theta_0$ is locally identified.  The set where the Jacobian drops rank is defined by the vanishing of a $(d_\theta \times d_\theta)$ minor of an infinite-dimensional Jacobian, which is a polynomial condition in $\theta$.  For a smooth, non-degenerate parametric family, this polynomial is not identically zero, so its zero set has Lebesgue measure zero.
\end{proof}

\begin{remark}[What generic identification does and does not say]\label{rem:generic}
Generic identification means: \emph{for almost every true parameter value, MLE is consistent}.  It does \emph{not} guarantee identification at every point---specific parameter configurations (e.g., at the boundary of the parameter space, or at isolated ``observationally equivalent'' points) may fail.  Nor does it speak to \emph{how well-identified} the model is: even if the model is identified, some parameters may be weakly identified if the Jacobian is close to singular, causing large standard errors in finite samples.
\end{remark}

\medskip\noindent
Generic identification is reassuring but imprecise.  We now verify identification for the specific parametric specification used in the simulation.


\subsection{Identification of the Beta Specification}\label{sec:beta_id}

We verify identification for the specification in the Monte Carlo:
\begin{align*}
    h \mid x &\sim \text{Beta}(a, b(x)) \text{ on } [h_{\min}, h_{\max}], \qquad b(x) = b_0 + b_1\, x
\end{align*}
with $x \in \{x_1, \ldots, x_{N_x}\}$ drawn from a discrete distribution (or, in the continuous case, from a parametric family $P(x; \psi)$).  In equilibrium, each $x_j$ maps to an entrant rate $r_2^*(x_j)$, inducing the conditional type distribution $F(h|r_2)$ and mixing distribution $G(r_2)$ as described in Remark~\ref{rem:induced}.  The full parameter vector is $\theta = (\alpha, \lambda, a, b_0, b_1, \psi)$, where $\psi$ parameterizes the distribution of~$x$.

\begin{prop}[Identification of the Beta specification]\label{prop:beta_id}
The parameter vector $\theta = (\alpha, \lambda, a, b_0, b_1, \psi)$ is identified from the joint distribution of $(r_i^{\text{win}}, W_i, D_i)$ if the following conditions hold:
\begin{enumerate}
    \item[(C1)] The support of $r_2$ among bank-2 winners contains at least three distinct values.
    \item[(C2)] $(a, b_0, b_1) \in (0,\infty) \times (0,\infty) \times \mathbb{R} \setminus \{0\}$ (the Beta shape parameter genuinely varies with $x$, and hence with $r_2$).
    \item[(C3)] $\alpha > 0$ and $\lambda \in \mathbb{R}$.
    \item[(C4)] $\Pr(W\!=\!2) > 0$ (some borrowers switch).
\end{enumerate}
\end{prop}

\begin{proof}
We proceed sequentially, using each data feature to pin down a subset of parameters.

\medskip\noindent
\textbf{Step 1: The functional restriction from $E[D|r_2, W\!=\!2]$.}

Define the function $\delta(r_2) \equiv E[D|r_2, W\!=\!2]$, which is identified from the data.  The model predicts:
\begin{align}\label{eq:delta}
    \delta(r_2) = \frac{\int_{h_{\min}}^{h_{\max}} h \cdot (1-s(\sigma^*(h;r_2),r_2)) \cdot f(h \mid r_2; \theta)\, dh}{\int_{h_{\min}}^{h_{\max}} (1-s(\sigma^*(h;r_2),r_2)) \cdot f(h \mid r_2; \theta)\, dh}
\end{align}
where $f(h|r_2;\theta)$ is the induced conditional type density at entrant rate~$r_2$ (inheriting the Beta shape from the primitive $F(h|x)$ through the equilibrium mapping), and $\sigma^*(h;r_2)$ depends on $(\alpha, \lambda)$.

This single functional equation in $r_2$ restricts three parameters: $(\alpha, \lambda, a, b_0, b_1)$ jointly.  The key insight is that these five parameters affect $\delta(r_2)$ through distinct channels:
\begin{itemize}
    \item $(\alpha, \lambda)$ control the \emph{selection weights} $(1-s)$: how much the adverse-selection reweighting distorts the mean type.  Specifically, $\alpha$ governs the sensitivity of the switching probability to rate differences, while $\lambda$ shifts the level.
    \item $(a, b_0, b_1)$ control the \emph{base distribution}: the unweighted mean $E[h|r_2]$ and its response to $r_2$.
\end{itemize}

By (C1), $\delta(r_2)$ is observed at three or more values of $r_2$.  Since $\delta(r_2)$ is not linear in any of the five parameters (due to the nonlinear equilibrium map $\sigma^*$ and the nonlinear Beta density), three or more values generically provide at least three independent nonlinear equations.  

Moreover, the \emph{shape} of $\delta(r_2)$ as a function of $r_2$ is particularly informative.  By (C2), $b_1 \neq 0$, so the Beta distribution genuinely shifts with $r_2$.  The curvature of $\delta(r_2)$ depends jointly on how the Beta distribution changes ($b_1$) and how the selection weights change ($\alpha$).  If $\alpha$ is large, selection is strong and $\delta$ is pulled further below $E[h|r_2]$; if $\alpha$ is small, $\delta \approx E[h|r_2]$ and the curvature is driven by the Beta family alone.

\medskip\noindent
\textbf{Step 2: The aggregate switching rate.}

The data identify $\Pr(W\!=\!2)$.  The model predicts:
\begin{align}\label{eq:agg_switch}
    \Pr(W\!=\!2) = \int \!\int (1 - s(\sigma^*(h;r_2), r_2)) \cdot f(h|r_2;\theta)\, dh\, dG(r_2;\theta)
\end{align}
This depends on $(\alpha, \lambda, \theta_F, \psi)$.  Crucially, $\lambda$ enters as an additive shifter in the logit: $s = 1/(1+\exp(\alpha\Delta - \lambda))$.  The aggregate switching rate is therefore monotonically decreasing in $\lambda$ (holding other parameters fixed), which means that once $(\alpha, \theta_F, \psi)$ are pinned down, $\lambda$ is uniquely determined by matching $\Pr(W\!=\!2)$.  This provides the primary direct restriction on $\lambda$.

\medskip\noindent
\textbf{Step 3: Separation of $(\alpha, \lambda)$ from $\theta_F$.}

The key separating mechanism is that $\alpha$ and $\lambda$ affect both the \emph{level} and the \emph{within-$r_2$ shape} of the default curve $\delta(r_2)$, while $\theta_F$ affects only the \emph{cross-$r_2$ pattern}.  Formally:
\begin{itemize}
    \item Fix $r_2$.  The selection weight $(1-s(\sigma^*(h;r_2),r_2))$ is a function of $h$ alone (via $\sigma^*$), parameterized by $(\alpha, \lambda)$.  Higher $\alpha$ makes this weight steeper in $h$ (stronger selection), pulling $\delta(r_2)$ further below $E[h|r_2]$.  The \emph{gap} $E[h|r_2] - \delta(r_2)$ at each $r_2$ disciplines $\alpha$.
    \item Across $r_2$ values: $\delta(r_2)$ traces a curve whose slope and curvature are governed by $b_1$ (how the type distribution shifts) and $a$ (its shape).  Since $b_1$ enters through the Beta family and $\alpha$ enters through the selection weights---and these are multiplicative in the integrand \eqref{eq:delta}---they generically cannot mimic each other.
\end{itemize}

\medskip\noindent
\textbf{Step 4: Identification of $\psi$.}

Given $(\alpha, \lambda, \theta_F)$, the selection probability $\Pr(W\!=\!2|r_2)$ and the equilibrium mapping $x \mapsto r_2^*(x)$ are both determined.  The induced population density of~$r_2$ must satisfy:
\begin{align*}
    g_{r_2|W=2}(r_2) = \frac{g(r_2;\theta) \cdot \Pr(W\!=\!2|r_2)}{\Pr(W\!=\!2)}
\end{align*}
Since $g_{r_2|W=2}$ is known from data and $\Pr(W\!=\!2|r_2)$ is now determined, $g(r_2;\theta) \propto g_{r_2|W=2}(r_2) / \Pr(W\!=\!2|r_2)$.  This pins down the population distribution of~$r_2$, and since $r_2^*(x)$ is monotone, it pins down the distribution of~$x$ and hence~$\psi$.

\medskip\noindent
\textbf{Step 5: Overidentification from bank-1 winners.}

Steps~1--4 use only bank-2 winner moments and the aggregate switching rate.  The entire distribution of $(r_1, D)$ among bank-1 winners, given by \eqref{eq:b1_density}, must also be matched by the same $\theta$.  This provides a rich set of overidentifying restrictions.  In particular, the bank-1 data independently constrain $(\alpha, \lambda)$ through the shape of the winning-rate distribution and the conditional default rates $E[D|r_1, W\!=\!1]$.  Any discrepancy between the parameter values implied by bank-1 vs.\ bank-2 moments signals misspecification.
\end{proof}

\begin{remark}[The role of $b_1 \neq 0$]\label{rem:b1}
Condition (C2) requires that the type distribution genuinely varies with $r_2$.  If $b_1 = 0$, the Beta distribution is the same at every $r_2$, and the function $\delta(r_2)$ varies only through the selection weights $(\alpha, \lambda)$.  In this case, $\delta(r_2)$ traces the composition of a fixed Beta distribution under varying selection, which restricts $(\alpha, \lambda, a, b_0)$ but reduces the number of independent restrictions.  Identification may still hold (from bank-1 moments), but it is less robust.
\end{remark}

\begin{remark}[Necessity of three $r_2$ values]\label{rem:three_r2}
Condition (C1) requires at least three distinct $r_2$ values among bank-2 winners.  With only one $r_2$ value, we are in the M5 setting and identification follows from the nonparametric argument (no parametric $F$ is needed).  With two $r_2$ values, $\delta(r_2)$ provides two restrictions, which combined with $\Pr(W\!=\!2)$ gives three equations for five parameters $(\alpha, \lambda, a, b_0, b_1)$---generically insufficient.  Three values provide the additional curvature information needed to separate the selection effect ($\alpha$) from the distribution shift ($b_1$).
\end{remark}

\begin{remark}[Practical verification via Monte Carlo]\label{rem:monte_carlo}
The propositions above are \emph{population} identification results: they state that the true parameter is the unique maximizer of the population log-likelihood.  Whether MLE can recover the parameters in finite samples depends on the strength of identification (the curvature of the likelihood near the optimum).  The Monte Carlo exercise in Section~\ref{sec:simulation} provides the practical verification: if MLE consistently recovers the true parameters across replications and parameter configurations, this confirms both identification and adequate finite-sample performance.
\end{remark}

\subsection{Intuition: What Each Data Feature Identifies}\label{sec:intuition}

\begin{center}
\renewcommand{\arraystretch}{1.5}
\begin{tabular}{p{3.5cm}p{9.5cm}}
\toprule
\textbf{Parameter(s)} & \textbf{Primary identifying variation} \\
\midrule
$\psi$ \newline (distribution of $\mathbf{x}$) & Distribution of winning rates among bank-2 winners, after correcting for selection.  Since $r_2^*(\mathbf{x})$ maps covariates to entrant rates, the population distribution of $r_2$ (and hence $\psi$) is recovered by inverting the selection distortion: $g(r_2;\theta) \propto g_{r_2|W=2}(r_2)/\Pr(W\!=\!2|r_2)$.  The selection correction depends on $(\alpha, \lambda, \theta_F)$, so $\psi$ is jointly identified with all other parameters. \\[4pt]
$\theta_F$ \newline (type distribution) & \textbf{How default rates vary with $r_2$ among bank-2 winners.}  If the default rate rises steeply as $r_2$ increases (riskier borrowers face higher entrant rates), this reveals how the type distribution shifts with $r_2$.  The shape of this curve disciplines $\theta_F$ through equation \eqref{eq:d2_restriction}. \\[4pt]
$\alpha$ \newline (price sensitivity) & \textbf{The severity of adverse selection.} $\alpha$ controls how much the selection function $1-s(h)$ varies with $h$.  High $\alpha$: strong selection, bank-2 winners are much better types than bank-1 winners.  Low $\alpha$: weak selection, default rates are similar across winners.  The \emph{gap} $E[D|W\!=\!2, r_2] - E[D|W\!=\!1, r_1]$ across matched rate levels traces out $\alpha$. \\[4pt]
$\lambda$ \newline (switching cost) & \textbf{The overall retention rate.}  $\lambda$ shifts bank-1's share uniformly: higher $\lambda$ means $\Pr(W\!=\!1)$ is higher at every rate differential.  The aggregate bank-1 share $\Pr(W\!=\!1)$ provides a direct moment, and conditional shares $\Pr(W\!=\!1|r_2)$ provide a functional restriction. \\
\bottomrule
\end{tabular}
\end{center}

\subsection{Conditions for Identification Failure}\label{sec:failure}

Even under Assumption~\ref{ass:parametric}, identification can fail in degenerate cases:

\begin{enumerate}
    \item \textbf{No bank-2 winners.} If $\lambda$ is very large, $\Pr(W\!=\!2) \approx 0$ and the data contain almost no bank-2 winners.  Without bank-2 winners, we lose the direct observation of $r_2^*$ and the conditional default rates that discipline $\theta_F$ and $\psi$.  In the limit $\Pr(W\!=\!2) = 0$, only bank-1 rates and defaults are observed, and the model collapses to a single-agent pricing problem where $(\alpha, \lambda)$ enter only through a single index $\lambda - \alpha\Delta$---a classic underidentification problem.
    
    \item \textbf{No variation in $r_2^*$.} If $\mathbf{x}$ has no effect on $r_2^*$ (or $\mathbf{x}$ is degenerate), the model reduces to M5 where bank~2 plays a single rate.  In this case, identification follows from the M5 argument and the parametric assumptions are unnecessary---but if we \emph{incorrectly} estimate a heterogeneous model, $\psi$ compresses to a point mass and $\theta_F$ effectively loses its $r_2$-dependence, which may cause numerical difficulties.
    
    \item \textbf{Overparameterization of $\theta_F$.} If the type distribution $F(h|\mathbf{x};\theta_F)$ has too many free parameters relative to the identifying variation (e.g., allowing flexible higher moments that the data cannot discipline), the model may be locally unidentified even if globally consistent.
\end{enumerate}

\subsection{Overidentifying Restrictions}\label{sec:overid}

Under the parametric specification, the model is typically \emph{over}identified: the data provide more restrictions than there are free parameters.  Key testable restrictions:

\begin{enumerate}
    \item \textbf{Consistency of bank-1 and bank-2 moments.}  The parameters $(\alpha, \lambda, \theta_F)$ are disciplined by bank-2 winners (through \eqref{eq:d2_restriction}) and by bank-1 winners (through \eqref{eq:b1_density}).  Any discrepancy between the parameter values implied by each subsample indicates model misspecification.
    
    \item \textbf{Adverse selection direction.}  The model predicts $E[D|W\!=\!2, r_2] > E[h|r_2]$ (bank-2 winners are riskier-than-average types, because $1-s$ is increasing in~$h$).  This is testable: if bank-2 winners have \emph{lower} default rates than bank-1 winners with similar rates, the adverse selection story is contradicted.
    
    \item \textbf{Linearity of the identifying restriction within strata.}  When conditioning on $r_2$ (observed for bank-2 winners), the M5 identifying restriction $\ln(\alpha m - 1) = \lambda - \alpha\Delta$ should hold.  Bank-2 winners reveal $r_2$, and if we could observe the bank-1 rate that each bank-2 winner \emph{would have} received, we could check this directly.  Instead, the restriction is embedded in the likelihood through $\sigma^*(h;r_2)$.
    
    \item \textbf{Bank-2's FOC.}  Having estimated $(\alpha, \lambda, \theta_F)$, the model predicts a specific $r_2^*(r_2)$ at each covariate value---i.e., the bank-2 FOC \eqref{eq:r2} should be satisfied.  The predicted distribution of $r_2^*$ must match the observed distribution among bank-2 winners (after selection correction).
\end{enumerate}


%===========================================================================
%  SECTION 5: LIKELIHOOD
%===========================================================================

\newpage
\section{Likelihood}\label{sec:likelihood}

The likelihood follows from the identification argument.

\subsection{Bank-1 Winners ($W_i = 1$)}

The entrant rate $r_2$ is latent; we marginalize over $G$:
\begin{align}\label{eq:L1}
    \boxed{\mathcal{L}_i^{(1)}(\theta) = \int_{\mathcal{R}_2(r_i^{\text{win}})} f\!\big(\tau(r_i^{\text{win}}; r_2) \mid r_2;\theta\big) \cdot \tau'(r_i^{\text{win}}; r_2) \cdot s(r_i^{\text{win}}, r_2) \cdot h^{D_i}(1\!-\!h)^{1-D_i} \cdot g(r_2;\theta)\, dr_2}
\end{align}
with $h = \tau(r_i^{\text{win}}; r_2)$.

\subsection{Bank-2 Winners ($W_i = 2$)}

The entrant rate is observed ($r_i^{\text{win}} = r_2$); the type $h$ is latent:
\begin{align}\label{eq:L2}
    \boxed{\mathcal{L}_i^{(2)}(\theta) = g(r_i^{\text{win}};\theta) \cdot \int_{\underline{h}}^{\bar{h}} \big(1-s(\sigma^*(h;r_i^{\text{win}}), r_i^{\text{win}})\big) \cdot h^{D_i}(1-h)^{1-D_i} \cdot f(h \mid r_i^{\text{win}};\theta)\, dh}
\end{align}

\subsection{Full Log-Likelihood}

\begin{align}\label{eq:L_full}
    \ell(\theta) = \sum_{i=1}^N \Big[\mathbf{1}(W_i=1)\ln \mathcal{L}_i^{(1)}(\theta) + \mathbf{1}(W_i=2)\ln \mathcal{L}_i^{(2)}(\theta)\Big]
\end{align}


\subsection{Computational Strategy}\label{sec:computation}

\begin{enumerate}
    \item Pre-compute $\sigma^*(h; r_2)$ on a fine grid of $(h, r_2)$ values.
    \item Construct $\tau(r_1; r_2)$ by numerical inversion on the same grid.
    \item Evaluate the integrals in \eqref{eq:L1}--\eqref{eq:L2} by quadrature.
    \item Key efficiency: the equilibrium grid depends only on $(\alpha, \lambda)$, not on $(\theta_F, \psi)$.  When only $\theta_F$ or $\psi$ changes, reuse the pre-computed grid.
\end{enumerate}


%===========================================================================
%  SECTION 6: MONTE CARLO SIMULATION
%===========================================================================

\newpage
\section{Monte Carlo Simulation}\label{sec:simulation}

The identification theory is validated through Monte Carlo simulation. The code is in \texttt{Code/model6/}.

\subsection{Data Generating Process}\label{sec:dgp}

\textbf{Grids.} $x \in \{x_1, \ldots, x_{N_x}\}$ and $h \in \{h_1, \ldots, h_{N_h}\}$ on equally spaced grids.

\textbf{Joint distribution.} $P(x=x_j, h=h_k)$ specified as an $N_x \times N_h$ probability matrix.  The conditional $P(h|x)$ is constructed using a Beta density with shape parameter depending on $x$, so that higher $x$ implies riskier borrowers.  Parameters are set in \texttt{params.m}.

\textbf{Equilibrium.} Solved for each $x_j$ using the nested algorithm (iterate on $r_2^*$ with damped best-response updates; within each iteration, solve bank-1's FOC by Brent's method for each $h_k$).

\textbf{Data generation.} $N$ borrowers drawn from $P(x,h)$; equilibrium rates assigned; winner drawn from logit shares; default drawn from Bernoulli($h$).

\subsection{Identification Checks}\label{sec:id_checks}

\begin{enumerate}
    \item \textbf{Bank-2 rate constancy within $x$.}  Among bank-2 winners with the same $x$, all winning rates should be identical ($r_i^{\text{win}} = r_2^*(x)$).  Confirms equilibrium is correctly computed.
    
    \item \textbf{Type recovery (conditional on $x$).}  Among bank-1 winners, $E[D_i \mid r_i^{\text{win}}, x_i]$ should equal $h_i$ (the M5 argument applies within each $x$-stratum).
    
    \item \textbf{Linear identifying restriction.}  $\ln(\alpha m - 1) = \lambda - \alpha\Delta$ should hold across all $(h, x)$ pairs.  Verified on population and simulated data.
    
    \item \textbf{Adverse selection.}  $E[D|W\!=\!2, x] > E[D|W\!=\!1, x]$ for each $x$ value.
\end{enumerate}

\subsection{Estimation Exercise}\label{sec:estimation_exercise}

The identification analysis is tested by:
\begin{enumerate}
    \item Simulating data under multiple parameter configurations (varying $\alpha$, $\lambda$, $\theta_F$).
    \item Estimating the model via MLE for each configuration.
    \item Comparing estimated to true parameters.
\end{enumerate}
The driver script \texttt{run\_identification.m} automates this, saving simulation and estimation results per configuration.  If the model is identified, the MLE should recover the true parameters as $N$ grows.


%===========================================================================
%  SECTION 7: DISCUSSION
%===========================================================================

\newpage
\section{Discussion}\label{sec:discussion}

\subsection{Comparison of Identification Across Models}

\begin{center}
\renewcommand{\arraystretch}{1.5}
\begin{tabular}{p{2.5cm}p{3.5cm}p{3.5cm}p{3.5cm}}
\toprule
& \textbf{M5} \newline (no heterogeneity) & \textbf{M6 with $\mathbf{x}$} \newline (covariates observed) & \textbf{M6 without $\mathbf{x}$} \newline (this document) \\
\midrule
$r_2$ observation & Single constant & Constant within $\mathbf{x}$ & Distribution \\[2pt]
Type recovery & Exact: \newline $E[D|r_1, W\!=\!1] = h$ & Exact within $\mathbf{x}$: \newline $E[D|r_1, W\!=\!1, \mathbf{x}] = h$ & Mixture: \newline $E[D|r_1, W\!=\!1]$ is a weighted avg \\[2pt]
$(\alpha, \lambda)$ & Nonparametric \newline (linear restriction) & Nonparametric \newline (overidentified) & Parametric \newline (via $F$, $G$ families) \\[2pt]
$F$ & Nonparametric & Nonparametric within $\mathbf{x}$ & Parametric $F(h|\mathbf{x};\theta_F)$ \\[2pt]
Key assumption & Logit demand & Logit demand & Logit + parametric $F$, $G$ \\
\bottomrule
\end{tabular}
\end{center}

\subsection{Sensitivity to Parametric Assumptions}

The parametric assumptions in Assumption~\ref{ass:parametric} are the price of not observing $\mathbf{x}$.  Several strategies can test sensitivity:

\begin{enumerate}
    \item \textbf{Flexible families.}  Use distributions with enough parameters to nest multiple shapes (e.g., 4-parameter beta, or a mixture of betas for $F$; a mixture of normals for $G$).  If estimates are stable across flexible specifications, the parametric assumption is not driving the results.
    
    \item \textbf{Partial observability.} When some covariates $\mathbf{z} \subseteq \mathbf{x}$ are observed, conditioning on $\mathbf{z}$ tightens the mixtures and reduces dependence on parametric assumptions.  Comparing estimates with and without conditioning on $\mathbf{z}$ reveals how much the parametric assumptions matter.
    
    \item \textbf{Specification tests.} The overidentifying restrictions in Section~\ref{sec:overid} provide formal tests of the parametric model.  Rejection indicates the parametric family is too restrictive.
\end{enumerate}

\subsection{Extensions}

\begin{enumerate}
    \item \textbf{Multi-product data.}  If the same borrower has multiple loan products, the common $(h, \mathbf{x})$ generates correlated rates across products, providing additional identification.
    
    \item \textbf{Panel data.}  Within-borrower variation in $r_2^*$ (driven by $\mathbf{x}$ changes over time) can be used to condition on $r_2$, reducing the mixture problem.
    
    \item \textbf{Semiparametric approaches.}  Rather than fully parametric $F$ and $G$, one could use sieve estimation or Bayesian nonparametric priors (e.g., Dirichlet process mixtures) to flexibly approximate the distributions while maintaining enough structure for identification.
\end{enumerate}


\end{document}
